\PassOptionsToPackage{unicode=true}{hyperref} % options for packages loaded elsewhere
\PassOptionsToPackage{hyphens}{url}
%
\documentclass[]{article}
\usepackage{lmodern}
\usepackage{amssymb,amsmath}
\usepackage{ifxetex,ifluatex}
\usepackage{fixltx2e} % provides \textsubscript
\ifnum 0\ifxetex 1\fi\ifluatex 1\fi=0 % if pdftex
  \usepackage[T1]{fontenc}
  \usepackage[utf8]{inputenc}
  \usepackage{textcomp} % provides euro and other symbols
\else % if luatex or xelatex
  \usepackage{unicode-math}
  \defaultfontfeatures{Ligatures=TeX,Scale=MatchLowercase}
\fi
% use upquote if available, for straight quotes in verbatim environments
\IfFileExists{upquote.sty}{\usepackage{upquote}}{}
% use microtype if available
\IfFileExists{microtype.sty}{%
\usepackage[]{microtype}
\UseMicrotypeSet[protrusion]{basicmath} % disable protrusion for tt fonts
}{}
\IfFileExists{parskip.sty}{%
\usepackage{parskip}
}{% else
\setlength{\parindent}{0pt}
\setlength{\parskip}{6pt plus 2pt minus 1pt}
}
\usepackage{hyperref}
\hypersetup{
            pdfborder={0 0 0},
            breaklinks=true}
\urlstyle{same}  % don't use monospace font for urls
\setlength{\emergencystretch}{3em}  % prevent overfull lines
\providecommand{\tightlist}{%
  \setlength{\itemsep}{0pt}\setlength{\parskip}{0pt}}
\setcounter{secnumdepth}{0}
% Redefines (sub)paragraphs to behave more like sections
\ifx\paragraph\undefined\else
\let\oldparagraph\paragraph
\renewcommand{\paragraph}[1]{\oldparagraph{#1}\mbox{}}
\fi
\ifx\subparagraph\undefined\else
\let\oldsubparagraph\subparagraph
\renewcommand{\subparagraph}[1]{\oldsubparagraph{#1}\mbox{}}
\fi

% set default figure placement to htbp
\makeatletter
\def\fps@figure{htbp}
\makeatother

\usepackage[]{natbib}
\bibliographystyle{plainnat}

\date{}

\begin{document}

\hypertarget{post-processing-of-fmriprep-outputs}{%
\subsubsection{Post-processing of fmriprep
outputs}\label{post-processing-of-fmriprep-outputs}}

The eXtensible Connectivity Pipeline- DCAN (XCP-D)
\citep{mehta2024xcp, mitigating_2018, satterthwaite_2013} was used to
post-process the outputs of \emph{fMRIPrep} version 22.0.0
\citep[RRID:SCR\_016216]{esteban2019fmriprep, esteban2020analysis}.
XCP-D was built with \emph{Nipype} version 1.10.0
\citep[RRID:SCR\_002502]{nipype1}.

\hypertarget{segmentations}{%
\paragraph{Segmentations}\label{segmentations}}

The following atlases were used in the workflow: the Schaefer
Supplemented with Subcortical Structures (4S) atlas
\citep{Schaefer_2017, pauli2018high, king2019functional, najdenovska2018vivo, glasser2013minimal}
at 10 different resolutions (1056, 156, 256, 356, 456, 556, 656, 756,
856, and 956 parcels), the Glasser atlas \citep{Glasser_2016}, the
Gordon atlas \citep{Gordon_2014}, the Tian subcortical atlas
\citep{tian2020topographic}, the HCP CIFTI subcortical atlas
\citep{glasser2013minimal}, the MIDB precision brain atlas derived from
ABCD data and thresholded at 75\% probability
\citep{hermosillo2024precision}, and the Myers-Labonte infant atlas
thresholded at 50\% probability \citep{myers2023functional}.

\hypertarget{anatomical-data}{%
\paragraph{Anatomical data}\label{anatomical-data}}

Native-space T1w images were transformed to MNI152NLin2009cAsym space at
1 mm3 resolution.

\hypertarget{functional-data}{%
\paragraph{Functional data}\label{functional-data}}

For each of the four BOLD runs found per subject (across all tasks and
sessions), the following post-processing was performed.

Non-steady-state volumes were extracted from the preprocessed confounds
and were discarded from both the BOLD data and nuisance regressors.
Framewise displacement was calculated from the motion parameters using
the formula from \citet{power_fd_dvars}, with a head radius of 50 mm.
Volumes with framewise displacement greater than 0.2 mm were flagged as
high-motion outliers for the sake of later censoring
\citep{power_fd_dvars}. In total, 36 nuisance regressors were selected
from the preprocessing confounds, according to the `36P' strategy. These
nuisance regressors included six motion parameters, mean global signal,
mean white matter signal, mean cerebrospinal fluid signal with their
temporal derivatives, and quadratic expansion of six motion parameters,
tissue signals and their temporal derivatives
\citep{benchmarkp, satterthwaite_2013}.

The BOLD data were converted to NIfTI format, despiked with
\emph{AFNI}'s \emph{3dDespike}, and converted back to CIFTI format.

Nuisance regressors were regressed from the BOLD data using a denoising
method based on \emph{Nilearn}'s approach. Any volumes censored earlier
in the workflow were first cubic spline interpolated in the BOLD data.
Outlier volumes at the beginning or end of the time series were replaced
with the closest low-motion volume's values, as cubic spline
interpolation can produce extreme extrapolations. The timeseries were
band-pass filtered using a(n) second-order Butterworth filter, in order
to retain signals between 0.01-0.08 Hz. The same filter was applied to
the confounds. The resulting time series were then denoised via linear
regression, in which the low-motion volumes from the BOLD time series
and confounds were used to calculate parameter estimates, and then the
interpolated time series were denoised using the low-motion parameter
estimates. The interpolated time series were then censored using the
temporal mask. The denoised BOLD was then smoothed using
\emph{Connectome Workbench} with a Gaussian kernel (FWHM=6 mm).

The amplitude of low-frequency fluctuation (ALFF) \citep{alff} was
computed by transforming the mean-centered, standard
deviation-normalized, denoised BOLD time series to the frequency domain
using the Lomb-Scargle periodogram
\citep{lomb1976least, scargle1982studies, townsend2010fast, taylorlomb}.
The power spectrum was computed within the 0.01-0.08 Hz frequency band
and the mean square root of the power spectrum was calculated at each
voxel to yield voxel-wise ALFF measures. The resulting ALFF values were
then multiplied by the standard deviation of the denoised BOLD time
series to retain the original scaling. The ALFF maps were smoothed with
the Connectome Workbench using a Gaussian kernel (FWHM=6 mm).

For each hemisphere, regional homogeneity (ReHo)
\citep{jiang2016regional} was computed using surface-based \emph{2dReHo}
\citep{surface_reho}. Specifically, for each vertex on the surface, the
Kendall's coefficient of concordance (KCC) was computed with
nearest-neighbor vertices to yield ReHo. For the subcortical, volumetric
data, ReHo was computed with neighborhood voxels using \emph{AFNI}'s
\emph{3dReHo} \citep{taylor2013fatcat}.

Processed functional timeseries were extracted from residual BOLD using
Connectome Workbench \citep{marcus2011informatics} for the atlases.
Corresponding pair-wise functional connectivity between all regions was
computed for each atlas, which was operationalized as the Pearson's
correlation of each parcel's unsmoothed timeseries with the Connectome
Workbench. In cases of partial coverage, uncovered vertices (values of
all zeros or NaNs) were either ignored (when the parcel had
\textgreater{}50.0\% coverage) or were set to zero (when the parcel had
\textless{}50.0\% coverage).

Many internal operations of \emph{XCP-D} use \emph{AFNI}
\citep{cox1996afni, cox1997software},\emph{Connectome Workbench}
\citep{marcus2011informatics}, \emph{ANTS} \citep{avants2009advanced},
\emph{TemplateFlow} version 25.1.1 \citep{ciric2022templateflow},
\emph{matplotlib} version 3.10.5 \citep{hunter2007matplotlib},
\emph{Nibabel} version 5.3.2 \citep{brett_matthew_2022_6658382},
\emph{Nilearn} version 0.13.0 \citep{abraham2014machine}, \emph{numpy}
version 2.2.6 \citep{harris2020array}, \emph{pybids} version 0.21.0
\citep{yarkoni2019pybids}, and \emph{scipy} version 1.15.3
\citep{2020SciPy-NMeth}. For more details, see the \emph{XCP-D} website
(https://xcp-d.readthedocs.io).

\hypertarget{copyright-waiver}{%
\paragraph{Copyright Waiver}\label{copyright-waiver}}

The above methods description text was automatically generated by
\emph{XCP-D} with the express intention that users should copy and paste
this text into their manuscripts \emph{unchanged}. It is released under
the \href{https://creativecommons.org/publicdomain/zero/1.0/}{CC0}
license.

\hypertarget{references}{%
\paragraph{References}\label{references}}

\bibliography{/usr/local/miniconda/lib/python3.10/site-packages/xcp\_d/data/boilerplate.bib}

\end{document}
